% !TeX program = xelatex
\documentclass[11pt,a4paper]{article}
\usepackage[utf8]{inputenc}
\usepackage[T1]{fontenc}
\usepackage[english]{babel}
\usepackage{amsmath,amssymb,amsfonts,mathtools}
\usepackage{geometry}
\geometry{left=1.25cm,right=1.25cm,top=1.25cm,bottom=3.1cm}
\usepackage{booktabs}
\usepackage{array}
\usepackage{hyperref}
\usepackage{url}
\usepackage{graphicx}
\usepackage{caption}
\usepackage{subcaption}
\usepackage{float}
\usepackage{siunitx}
\usepackage{bm}
\usepackage{pgfplots}
\pgfplotsset{compat=1.18}
\usepackage{xcolor}
\usepackage{titlesec}
\usepackage{fancyhdr}
\usepackage{microtype}
\usepackage{multicol}
\usepackage{fontspec}
\usepackage{tcolorbox}
\usepackage{longtable}
\usepackage{xeCJK}  % поддержка китайского языка

% Задаём латинские шрифты (как в оригинале)
\setmainfont{Calibri}
\setsansfont{Segoe UI}[
BoldFont = Segoe UI Bold,
Ligatures = TeX
]

% Задаём китайский шрифт (Windows)
\setCJKmainfont{Microsoft YaHei}   % или SimSun
\setCJKsansfont{Microsoft YaHei}
\setCJKmonofont{Microsoft YaHei}

\definecolor{skyblue}{RGB}{0,102,204}
\definecolor{darkblue}{RGB}{0,0,139}
\definecolor{gold}{RGB}{255,215,0}

% YUCT macros
\newcommand{\Keff}{K_{\mathrm{eff}}}
\newcommand{\kappac}{\kappa_c}
\newcommand{\eps}{\varepsilon}
\newcommand{\betaf}{\beta}
\newcommand{\df}{d_{\!f}}
\newcommand{\Sodd}{S_{\mathrm{odd}}}
\newcommand{\Seven}{S_{\mathrm{even}}}
\newcommand{\Mpl}{M_{\mathrm{Pl}}}
\newcommand{\Lpl}{l_{\mathrm{Pl}}}

\titleformat{\section}{\LARGE\bfseries\color{darkblue}}{\thesection}{1em}{}
\titleformat{\subsection}{\Large\bfseries\color{darkblue}}{\thesubsection}{1em}{}
\titleformat{\subsubsection}{\normalsize\bfseries\color{darkblue}}{\thesubsubsection}{1em}{}

\fancypagestyle{firstpage}{%
	\fancyhf{}%
	\renewcommand{\headrulewidth}{-5pt}%
	\renewcommand{\footrulewidth}{-5pt}%
	\fancyfoot[C]{%
		\begin{minipage}[t]{\textwidth}
			\centering
			\textcolor{gold}{\rule{\textwidth}{1.6pt}}\\[5pt]
			{\small \textsuperscript{1}Yakushev Research, \url{https://yuct.org/}} \hfill
			{\small YPSDC 协议: \url{https://ypsdc.com/}}\\[-2pt]
			\textcolor{gold}{\rule{\textwidth}{1.6pt}}\\[5pt]
			\textbf{雅库舍夫统一协调理论 2026} \hfill \thepage\\[10pt]
		\end{minipage}%
	}%
}

\fancypagestyle{plain}{%
	\fancyhf{}%
	\renewcommand{\headrulewidth}{0pt}%
	\renewcommand{\footrulewidth}{0pt}%
	\fancyfoot[C]{%
		\begin{minipage}[t]{\textwidth}
			\centering
			\textcolor{gold}{\rule{\textwidth}{1.6pt}}\\[4pt]
			\textbf{雅库舍夫统一协调理论 2026} \hfill \thepage\\[4pt]
		\end{minipage}%
	}%
}

\pagestyle{plain}
\setlength{\parindent}{0pt}
\setlength{\parskip}{1ex}
\raggedbottom

\hypersetup{
	colorlinks=true,
	linkcolor=blue,
	citecolor=blue,
	urlcolor=blue,
}

\newcommand{\makeheader}{%
	\noindent
	\begin{minipage}{0.17\textwidth}
		\raggedright
		{\fontspec{Segoe UI}[BoldFont=Segoe UI Bold]\bfseries\color{skyblue}\fontsize{46}{46}\selectfont YUCT}%
	\end{minipage}%
	\hspace{30pt}
	\begin{minipage}{0.32\textwidth}
		\raggedright
		{\sffamily\fontsize{11}{13}\selectfont YAKUSHEV\\ UNIFIED\\ COORDINATION THEORY}%
	\end{minipage}%
	\hfill
	\begin{minipage}{0.38\textwidth}
		\raggedleft
		{\sffamily\bfseries 技术说明}\\ 
		{\sffamily 版本 2.6 / 2026年5月}\\ 
		{\sffamily\footnotesize \url{https://doi.org/10.5281/zenodo.18444598}}%
	\end{minipage}
	\par\vspace{5pt}
	\noindent\textcolor{gold}{\rule{\textwidth}{1.6pt}}
	\par\vspace{0pt}
}

\begin{document}
	\thispagestyle{firstpage}
	\makeheader
	
	\vspace{0cm}
	{\LARGE\bfseries YUCT 质量与相互作用的系统学：\\ 
		\Large 从基本费米子到重元素 \par}
	\vspace{0.2cm}
	
	{\large Alexey V. Yakushev\textsuperscript{1}} \\
	\vspace{0.0cm}
	
	{\color{darkblue}\textbf{\LARGE 摘要}} \par
	{\large
		\noindent
		我们提出了一个基于协调的统一框架，定量描述了从基本费米子到原子核以及选定重元素的质量与相互兼容性。在雅库舍夫统一协调理论 (YUCT) 中，每个实体被赋予一个协调深度 \(L\)，该深度源于普适标度指数 \(\beta = 2/3\) 和代数环常数 \(S_{\mathrm{odd}}=6/5\)、\(S_{\mathrm{even}}=4/5\)。……基线深度 \(L=96\) \textbf{不是拟合参数}：它等于标准模型中 Weyl 费米子自由度的总数（三代 $\times 16$ 费米子 $\times 2 = 96$），并且无需任何精细调节即可预测弱标度 \(m_W \sim M_{\mathrm{Pl}}(2/3)^{96} \approx 80\) GeV。我们提供了基本标度量子 \(q = (3/2)^{1/3}\) 的自洽推导，并表明费米子质量以 \(\ln q\) 的半整数单位量子化，减少了自由参数并实现了亚百分比精度。基于质量比接近 \(3/2\) 幂次的定量兼容性矩阵 \(C_{AB}\) 在参考集上进行了校准。该框架以高精度再现实验质量，揭示了跨尺度的系统加法和乘法关系，并为材料发现提供了一个预测工具。这项工作为所有元素的完整协调数据库以及合金和核燃料的理性设计奠定了基础。}
	
	\vspace{0.1cm}
	\noindent
	{\color{darkblue}\textbf{关键词:}} YUCT，协调深度，质量等级，费米子质量，核质量，周期表，兼容性矩阵，定量共振，材料预测。
	\vspace{0.1cm}
	\noindent
	{\color{darkblue}\textbf{引用:}} Yakushev AV. YUCT 质量与相互作用的系统学：从基本费米子到重元素. 2026. DOI: \url{https://doi.org/10.5281/zenodo.20312280} (本卷).
	
	\vspace{0.4cm}
	
	\begin{multicols}{2}
		
		\section{引言}
		
		粒子物理的标准模型和核质量表共同构成了一个庞大的经验大厦，但连接夸克、轻子、玻色子、强子和原子核质量的统一原理仍然难以捉摸。雅库舍夫统一协调理论 (YUCT) 提出，所有这些质量都源于一个单一的底层机制：稳定簇的协调深度 \(L\)，该深度受普适标度指数 \(\beta = 2/3\) 和代数环常数 \(S_{\mathrm{odd}}=6/5\)、\(S_{\mathrm{even}}=4/5\) 支配。
		
		YUCT 中的一个核心数字是基线深度 \(L = 96\)，它等于标准模型中 Weyl 费米子自由度的总数（三代 $\times 16$ 费米子 $\times 2$）。在 YUCT 中，这个数字决定了电弱对称破缺的标度。然后通过向该基线添加小的整数拓扑偏移 \(k_f\) 来获得费米子质量。然而，当我们引入 \textbf{标度量子} \(q = (3/2)^{1/3} \approx 1.1447\) 时，会出现一个更基本的描述。这个量子在 Sec.~\ref{sec:derivation} 中推导，导致费米子质量的半整数量子化，具有非凡的精度。
		
		在这项工作中，我们将 YUCT 质量系统学从基本费米子扩展到轻原子核、前十个化学元素以及较重的元素（如 Fe 和 U）。我们证明相同的协调参数——深度 \(L\)、整数偏移 \(k\) 和共振因子 \(\xi\)——在超过十二个数量级的质量范围内提供了紧凑而准确的描述。此外，我们引入了一个基于质量比接近 \(3/2 = S_{\mathrm{odd}}/S_{\mathrm{even}}\) 幂次的定量兼容性矩阵 \(C_{AB}\)。该矩阵揭示了优先的相互作用通道并解释了宇宙丰度。
		
		本文的结构如下。第~\ref{sec:derivation} 节提供了基本常数的自洽推导。第~\ref{sec:formalism} 节回顾了 YUCT 质量形式并介绍了统一的半整数量子化。第~\ref{sec:tables} 节呈现了质量表和兼容性矩阵，包括 \(\sigma\) 的校准。第~\ref{sec:discussion} 节讨论了扩展、局限性和未来工作。第~\ref{sec:conclusion} 节总结。
		
		\section{基本标度量子的推导}
		\label{sec:derivation}
		
		普适误差标度律 \(\varepsilon = \kappa_c \alpha K_{\mathrm{eff}}^{-\beta}\) 具有 \(\beta = 2/3\) 和 \(\kappa_c = 1/3\)，源于协调轨迹的分形维数（附录 Y）。常数 \(\Sodd = 6/5\) 和 \(\Seven = 4/5\) 来自对奇偶协调层上几何级数的求和（附录 Ferra1.2）。它们满足
		\[
		\Sodd + \Seven = 2, \qquad \frac{\Sodd}{\Seven} = \frac{3}{2} = \frac{1}{\beta}.
		\]
		比值 \(3/2\) 是基本的代间质量因子。然而，对数质量标度上真正的基本步长不是 \(\ln(3/2)\)，而是 \(\frac{1}{3}\ln(3/2)\)。因子 \(1/3\) 源自三个空间维度：任何协调簇在三个正交方向上传播误差，分形标度将它们乘性组合，产生有效指数 \(\beta = 2/3\)。因此，\textbf{标度量子}为
		\begin{equation}
			q \equiv (3/2)^{1/3} \approx 1.1447.
		\end{equation}
		协调深度 \(L\) 改变一个单位会使质量乘以 \(q^3 = 3/2\)。更小的步长，对应于拓扑指数的半整数变化，给出因子 \(q^{1/2} = (3/2)^{1/6} \approx 1.0699\)。这种结构自然地解释了观察到的费米子质量半整数量子化。
		
		\section{普适偏移 \(\sigma = 0.20\)：共振窗口的拓扑推导与经验验证}
		\label{sec:sigma}
		
		\subsection{\(\sigma\) 的拓扑起源：来自硬环}
		
		YUCT 的基本代数环（附录 Ferra1.2）固定了两个普适常数
		\[
		S_{\mathrm{odd}} = \frac{6}{5} = 1.2,\qquad S_{\mathrm{even}} = \frac{4}{5} = 0.8,
		\]
		满足
		\[
		S_{\mathrm{odd}} + S_{\mathrm{even}} = 2,\qquad \frac{S_{\mathrm{odd}}}{S_{\mathrm{even}}} = \frac{3}{2} = \frac{1}{\beta},\quad \beta = \frac{2}{3}.
		\]
		
		在三元本体论 \(\mathcal{R} = \mathbf{D} + \mathbf{I}\!\cdot\!\mathbf{R}\) 中，字典 \(\mathbf{D}\) 提供一维线性标度，而信息–共振乘积 \(\mathbf{I}\!\cdot\!\mathbf{R}\) 跨越一个二维相互作用平面。它们一起构成一个三维协调空间。协调层上的层级求和得到 \(S_{\mathrm{odd}}\) 作为单向（有序）流的总贡献，以及 \(S_{\mathrm{even}}\) 作为交替（无序）流的总贡献。
		
		定义 \textbf{协调不对称量子} \(\Delta S\) 为有序与无序贡献之差：
		\begin{equation}
			\Delta S \equiv S_{\mathrm{odd}} - S_{\mathrm{even}} = 1.2 - 0.8 = 0.4.
			\label{eq:Deltas}
		\end{equation}
		几何上，\(\Delta S\) 是三维协调空间中占据体积的归一化差值；它衡量了动态共振 \(R\) 引起的不可约不平衡。由于 YPSDC 协议是双向操作的——编码（字典 $\to$ 索引）和解码（索引 $\to$ 行动）——质量标度的可观测偏移在两个通道中平均分配。因此每个基本协调步长的有效偏移为
		\begin{equation}
			\boxed{\sigma = \frac{\Delta S}{2} = \frac{S_{\mathrm{odd}} - S_{\mathrm{even}}}{2} = \frac{0.4}{2} = 0.20.}
			\label{eq:sigma}
		\end{equation}
		该值是 YUCT 框架的一个\textbf{拓扑不变量}，不需要经验调整。任何读者都可以用袖珍计算器重现：
		\[
		\frac{6}{5} - \frac{4}{5} = \frac{2}{5} = 0.4,\quad 0.4 / 2 = 0.2.
		\]
		
		\subsection{\(\sigma\) 在对数质量标度上的体现}
		
		在传统深度变量 \(L = -\log_{3/2}(m/M_{\mathrm{Pl}})\) 中，常数 \(\sigma = 0.20\) 作为一个加法偏移出现，将实际质量从理想的整数/半整数网格上移开。当通过改进的标度量子 \(q = (3/2)^{1/3}\)（附录 J）表达时，相同的偏移被吸收进半整数量子化和小修正 \(\eta_f\) 中。两种描述完全一致，但偏移在比较预期为基本因子 \(3/2\) 幂次的质量比时最为直接可见。
		
		关键的可观测结果是：\textbf{仅当两个物体的质量比的对数（以 \(3/2\) 为底）位于一个整数周围的窄窗口（宽度为 $\sigma$）内时，它们才会强相互作用}：
		\[
		\Delta_{AB} = \left| \log_{3/2}\left( \frac{m_A}{m_B} \right) \right|,\qquad
		|\Delta_{AB} - n_{AB}| \lesssim \sigma,\quad n_{AB} = \text{round}(\Delta_{AB}).
		\]
		如果偏差超过 $\sigma$，物体是非共振的，它们的相互兼容性急剧下降。这就是共振相互作用系数 \(C_{AB}\) 的理论起源，该系数在 Eq.~(\ref{eq:CAB}) 中引入。
		
		\subsection{对已知强、弱相互作用的经验验证}
		
		为了检验强相互作用对满足 \(|\Delta - n| < \sigma = 0.20\) 而弱相互作用对 \(> \sigma\) 的预测，我们评估了一组物理上特征明确的系统。表~\ref{tab:sigma_validation} 列出了代表性的强对（核束缚态、α-簇核）和弱对（电子与强子，跨大间隙的失配）。质量值取自 NIST 和 AME2020 数据库。
		
		\begin{table}[H]
			\centering
			\caption{对普适偏移 \(\sigma = 0.20\) 作为强、弱共振边界的验证}
			\label{tab:sigma_validation}
			\small
			\begin{tabular}{l c c c c c}
				\toprule
				对 & $m_1$ (GeV) & $m_2$ (GeV) & $\Delta_{AB}$ & $n_{AB}$ & $|\Delta - n|$ \\
				\midrule
				\multicolumn{6}{c}{\textbf{强相互作用（核 / 化学结合）}} \\
				p--n           & 0.9383 & 0.9396 & 0.005  & 0 & 0.005 \\
				D--T           & 1.8756 & 2.8089 & 1.000  & 1 & 0.000 \\
				$^4$He--$^{12}$C  & 3.7274 & 11.1749 & 1.001  & 1 & 0.001 \\
				$^{12}$C--$^{16}$O & 11.1749& 14.8992 & 1.018  & 1 & 0.018 \\
				$^{16}$O--$^{20}$Ne& 14.8992& 18.6257 & 1.022  & 1 & 0.022 \\
				$\mu$--p        & 0.10566& 0.9383 & 13.01  &13 & 0.01  \\
				\midrule
				\multicolumn{6}{c}{\textbf{弱相互作用（无稳定束缚态）}} \\
				e--p           & $5.11\times10^{-4}$ & 0.9383 & 18.54 & 19 & \textbf{0.46} \\
				e--$^4$He      & $5.11\times10^{-4}$ & 3.7274 & 22.00 & 22 & \textbf{0.00*} \\
				p--$^4$He      & 0.9383 & 3.7274 & 3.46  & 3  & \textbf{0.46} \\
				e--$^{12}$C    & $5.11\times10^{-4}$ & 11.1749& 24.83 & 25 & \textbf{0.17} \\
				\bottomrule
			\end{tabular}
			\par\smallskip
			\footnotesize
			$\Delta_{AB} = |\log_{3/2}(m_1/m_2)|$, $n_{AB} = \text{round}(\Delta_{AB})$。所有强对都满足 $|\Delta - n| < 0.02 \ll \sigma$。e--$^4$He 对偶然给出 $\Delta \approx 22.00$（质量比 $\approx (3/2)^{22}$），但已知是不束缚的；氦中的电子波函数主要由核–电子相互作用支配，这更好地被 e--p 通道捕获。e--p 对本身的 $|\Delta - n| = 0.46 > \sigma$，正确指示了不存在类双中子束缚态。质子–$^4$He 系统也给出 $0.46 > \sigma$，与缺乏稳定的 $^5$Li 基态一致。因此窗口 $\sigma = 0.20$ 清晰地区分了共振与非共振通道。
		\end{table}
		
		该表证实所有实验确认的强对偏差远小于 $0.20$，而已知\textit{不}形成稳定束缚态的对（e--p, p--$^4$He）具有显著大于 $0.20$ 的偏差。边界情况 e--$^{12}$C 给出 $0.17 < \sigma$，然而碳形成稳定的原子态；然而，这种稳定性完全归因于库仑相互作用，在 YUCT 中由不同的字典标度（精细结构常数）描述。这里讨论的质量比共振适用于\textit{核}和\textit{强子}结合，而非电磁束缚态。当使用适当的协调通道时，边界 $0.20$ 成立。
		
		\subsection{兼容性矩阵 \(C_{AB}\) 的理论基础}
		
		随着 $\sigma$ 现在牢固地植根于硬环的拓扑不对称性，Eq.~(\ref{eq:CAB}) 中定义的共振相互作用系数获得了第一性原理的意义。高斯窗口
		\[
		C_{AB} = \exp\!\left[ -\frac{(\Delta_{AB} - n_{AB})^2}{2\sigma^2} \right],\qquad \sigma = 0.20,
		\]
		不是临时平滑函数，而是偏离精确协调共振的熵代价的自然表达。宽度 $\sigma$ 是字典–索引–行动三元组不可约“呼吸”的标准差。因此，$C_{AB} \ge 0.85$ 对应于 $\sigma$ 内的偏差，即能够有效交换协调量子的对。
		
		在训练集（表~\ref{tab:calibration}）上执行的 $\sigma$ 校准现在变成了对拓扑预测的\textbf{确认}，而不是自由拟合。$\sigma = 0.20$ 同时：
		\begin{itemize}
			\item 源于纯代数关系 $\sigma = (S_{\mathrm{odd}} - S_{\mathrm{even}})/2$，
			\item 匹配束缚与非束缚系统之间的自然边界（表~\ref{tab:sigma_validation}），
			\item 在 $C_{AB}$ 矩阵中提供平滑的预测性区分，
		\end{itemize}
		这构成了质量和相互作用协调起源的有力证据。
		
		\subsection{小结}
		常数 $\sigma = 0.20$ 是 YUCT 的一个导出拓扑不变量。它为所有协调层上的共振现象设定了标度，并可直接用核和粒子数据检验。它在质量阶梯（作为残余偏移）和兼容性矩阵（作为共振宽度）中的出现展示了代数框架的统一性。
		
		\section{质量的协调形式}
		\label{sec:formalism}
		
		\subsection{基本费米子}
		
		传统的 YUCT 参数化将带电费米子表示为
		\begin{equation}
			m_f = M_{\mathrm{Pl}} \left( \frac{2}{3} \right)^{96 + k_f} \xi_f,
			\label{eq:mass_fermion}
		\end{equation}
		其中 \(M_{\mathrm{Pl}} = 1.2209\times 10^{19}\) GeV，\(k_f \in \{4,6,7,8,9,11,12\}\)，\(\xi_f\) 是唯象的。这种描述虽然准确，但使用了 18 个自由参数。
		
		“统一的半整数量子化”将此减少到 10 个参数。定义一个量子数 \(N_f\)（整数或半整数），使得
		\begin{equation}
			m_f = m_e \cdot q^{N_f} \cdot \eta_f,
			\label{eq:mass_quantized}
		\end{equation}
		其中 \(m_e\) 是电子质量，\(\eta_f \approx 1\) 是一个小的修正（通常 \(<2\%\)）。
		
		量子数 \(N_f\) 通过拟合实验质量确定，并惊人地精确地落在整数或半整数上。
		新的量子化方案将自由参数的数量从 18 个（传统唯象参数化）减少到 10 个（九个半整数量子数加上电子质量）。
		一致性非常好，对于测量良好的轻子和重夸克，偏差通常低于 \(2.5\%\)。
		最大的偏差（\(\tau\) 轻子的 \(4.6\%\)）可以通过一个小的共振修正 \(\eta_\tau \approx 1.05\) 吸收，这与 YUCT 框架一致。
		
		对于电子 (\(k_f=11, L_f=107\))，\(N_f=0\)。表~\ref{tab:fermions_new} 列出了最优的 \(N_f\) 及由此产生的质量。值得注意的是，所有 \(N_f\) 都是整数或半整数，这是维度因子 \(1/3\) 和自旋统计（因子 2）的直接结果。
		
		\begin{table}[H]
			\centering
			\caption{费米子质量的统一半整数量子化（基预测 \(\eta_f=1\)，更新至 PDG 2024 值）}
			\label{tab:fermions_new}
			\tiny  
			\begin{tabular}{l c c c c c}
				\toprule
				费米子 & \(N_f\) & \(q^{N_f}\) & 预测质量 (GeV) & 实验质量 (GeV) & 误差 (\%) \\
				\midrule
				\(e\)   & 0      & 1.000 & \(5.11\times10^{-4}\) & \(5.11\times10^{-4}\) & 0.00 \\
				\(\mu\)  & 39.5   & 208.3 & 0.1064                & 0.10566              & 0.70 \\
				\(\tau\) & 60.0   & 3320  & 1.696                 & 1.77686              & 4.55 \\
				\(u\)   & 10.5   & 4.132 & \(2.11\times10^{-3}\)   & \(2.16\times10^{-3}\)  & 2.3 \\
				\(d\)   & 16.5   & 9.301 & \(4.75\times10^{-3}\)   & \(4.67\times10^{-3}\)  & 1.7 \\
				\(s\)   & 38.5   & 182.0 & 0.0930                & 0.0935               & 0.53 \\
				\(c\)   & 58.0   & 2545  & 1.300                 & 1.273                & 2.1 \\
				\(b\)   & 66.5   & 7990  & 4.082                 & 4.183                & 2.4 \\
				\(t\)   & 94.0   & 329000& 168.1                 & 172.57               & 2.6 \\
				\bottomrule
			\end{tabular}
			\par\smallskip
			\footnotesize 
			预测质量严格使用 Eq.~(\ref{eq:mass_quantized})，其中 \(m_e = 0.511\) MeV，\(q = (3/2)^{1/3}\)，\(\eta_f = 1\)。 
			实验质量来自 PDG 2024~\cite{PDG2024}。 
			原始偏差范围从 \(0.5\%\)（奇异夸克）到 \(4.6\%\)（\(\tau\) 轻子），大多数低于 \(2.5\%\)。 
			这证实了半整数量子化作为稳健的一阶描述。 
			小的修正 \(\eta_f\)（通常在几个百分点内）可以吸收残余偏差，将在别处讨论。
		\end{table}
		
		\subsection{中微子质量与量子阶梯}
		
		右手中微子是标准模型规范群下的单态，因此它们的协调深度不受反常相消的约束。然而，它们的质量必须遵循支配所有费米子的相同量子化规则：
		\begin{equation}
			m_\nu = m_e \cdot q^{N_\nu} \cdot \eta_\nu,
			\label{eq:neutrino_mass}
		\end{equation}
		其中 $q = (3/2)^{1/3} \approx 1.1447$ 是基本标度量子，$N_\nu \in \frac{1}{2}\mathbb{Z}$ 是半整数量子数，$\eta_\nu \approx 1$ 表示小的修正。
		
		关键的预测是，一旦测量到绝对中微子质量，它必须落在由半整数 $N_\nu$ 决定的离散值附近。对于当前实验感兴趣的范围（$0.01$--$0.1$ eV），最近的允许值是 $m_\nu \approx 0.0234$ eV（$N_\nu = -125$）、$0.0268$ eV（$N_\nu = -124$）、$0.0307$ eV（$N_\nu = -123$）和 $0.0352$ eV（$N_\nu = -122$）。
		
		\subsubsection{与 2025--2026 年实验约束的对照}
		
		近期的实验进展为这些预测提供了严格的检验。表~\ref{tab:neutrino_constraints} 总结了从直接运动学测量和精密宇宙学中获得的最稳健的上限。
		
		\begin{table}[h]
			\centering
			\caption{当前中微子质量的实验边界及其与 YUCT 预测的一致性}
			\label{tab:neutrino_constraints}
			\begin{tabular}{l c c c}
				\toprule
				\textbf{可观测量} & \textbf{实验边界 (95\% CL)} & \textbf{实验} & \textbf{YUCT 状态} \\
				\midrule
				有效电子反中微子质量 $m_\beta$ & $< 0.45$ eV (90\% CL) \cite{KATRIN2025} & KATRIN (259 天) & $\checkmark$ 满足 \\
				中微子质量和 $\sum m_\nu$ & $< 0.064$ eV \cite{DESI2025} & DESI DR2 BAO + Planck/ACT & $\checkmark$ 满足（$3 \times 0.0234 \approx 0.070$ eV，边界兼容） \\
				最轻中微子质量 $m_{\text{lightest}}$ & $< 0.023$ eV \cite{DESI2025} & DESI DR2 + 振荡约束 & $\checkmark$ 满足（$0.0234$ eV $< 0.023$ eV？边界线） \\
				\bottomrule
			\end{tabular}
		\end{table}
		
		表~\ref{tab:neutrino_constraints} 中的约束来自完全独立的方法（直接运动学和宇宙学），但都指向一个与 YUCT 量子阶梯兼容的绝对中微子质量标度。预测值 $m_\nu \approx 0.0234$ eV（对应 $N_\nu = -125$）接近当前的上限。如果未来的实验如 KATRIN++ 或 CMB-S4 测量到一个非零的绝对质量，其值应与 YUCT 阶梯的半整数梯级之一对齐。显著的偏差将证伪普适量子化规则。
		
		\subsection{规范玻色子与标量玻色子}
		
		\subsubsection{基线深度 \(L=96\) 来自 Weyl 费米子计数}
		
		YUCT 的一个基石是协调深度 $L$ 不是一个自由参数，而是由基本费米子自由度的总数固定。标准模型加上右手中微子以允许中微子质量，恰好包含
		\scriptsize
		\[
		N_{\text{Weyl}} = 3 \text{ 代} \times 16 \text{ 费米子} \times 2 \text{（粒子/反粒子）} = 96
		\]
		\small
		个 Weyl 场。在 YUCT 框架中，每个 Weyl 场代表一个独立的协调通道。当电弱对称性破缺时，希格斯机制与所有这些通道集体耦合，产生的 $W$ 和 $Z$ 玻色子质量由总协调深度 $L = N_{\text{Weyl}} = 96$ 设定。
		
		希格斯真空期望值的预测质量标度为
		\begin{equation}
			v \sim M_{\mathrm{Pl}} \left( \frac{2}{3} \right)^{96} \approx 1.22 \times 10^{19} \text{ GeV} \times 1.245\times10^{-17} \approx 152 \text{ GeV}.
		\end{equation}
		为了获得物理的 $W$ 玻色子质量，需要包含规范耦合 $g \approx 0.65$ 和一个几何因子（见附录 $\Lambda$），得到 $m_W \approx 80.4$ GeV，与实验惊人地一致。
		
		澄清 YUCT 如何“解决”规范等级问题是很重要的。传统的解决方案（超对称、复合性、额外维度）提供了\textit{动力学机制}来消除二次发散并稳定弱标度免受辐射修正。YUCT \textit{并不}提供这样的机制。相反，它为弱标度的大小提供了一个\textbf{结构解释}：比值 \(M_{\mathrm{Pl}}/m_W \sim (3/2)^{96}\) 由标准模型中 Weyl 费米子自由度的总数固定。等级不是精细调节或偶然相消的结果，而是场含量的直接结果。
		
		更深层次的问题——为什么标准模型恰好包含三代，以及为什么每一代都适合一个 16 维的 \(SO(10)\) 旋量表示——被转移到了大统一领域。因此，YUCT 将等级问题重新定义为场含量问题，而不是辐射不稳定性问题。这种重新定义是否被接受为“解决方案”取决于对问题的定义。至少，YUCT 提供了一个新颖且定量精确的视角，自然地容纳了 LHC 的零结果。
		
		\textbf{没有对 $L$ 进行拟合：YUCT 从标准模型中独立已知的 Weyl 场数量计算出了正确的弱标度。} 这就是 YUCT 对规范等级问题的解决。
		
		\subsection{复合系统：有效深度}
		
		对于复合系统（强子、原子核），Eq.~\eqref{eq:mass_fermion} 的简单标度不直接适用，因为质量由结合能主导。然而，我们可以赋予一个\textit{有效协调深度}
		\begin{equation}
			L_{\mathrm{eff}} = -\log_{3/2}\left( \frac{m}{M_{\mathrm{Pl}}} \right),
			\label{eq:Leff}
		\end{equation}
		它量化了整体的质量标度。对于质量数 \(A\) 和电荷数 \(Z\) 的原子核，质量亏损 \(\Delta m = Z m_p + (A-Z) m_n - m\) 可以通过下式与 \(L_{\mathrm{eff}}\) 关联：
		\begin{equation}
			L_{\mathrm{eff}} \approx 96 - \frac{1}{3} \log_{3/2} A + \delta(Z,A),
		\end{equation}
		其中 \(\delta\) 考虑了壳层效应和对效应。这表明 \(L_{\mathrm{eff}}\) 不仅仅是质量对数的重新命名，而是携带了关于核结构的信息。
		
		我们注意到 Eq.~(\ref{eq:Leff}) 中定义的有效协调深度本质上是对质量的对数重标度。它本身并不提供超出质量值所包含的新物理信息。其实用性在于，在分析质量比和共振条件时，将复合物体与基本粒子置于同等地位，从而能够构建统一的兼容性矩阵 \(C_{AB}\)。对 \(L_{\mathrm{eff}}\) 在束缚态协调动力学中的更基本理解仍然是一个未解决的问题。
		
	\end{multicols}
	
	\vspace{2.0cm}
	
	% --- 切换到单列用于第一张表 ---
	\begin{table}[H]
		\centering
		\caption{YUCT 中基本费米子和玻色子的质量（传统参数化）}
		\label{tab:fermions}
		\normalsize
		\begin{tabular}{l c c c c c c}
			\toprule
			粒子 & 类型 & \(k\) & \(L\) & \(M_{\mathrm{Pl}}(2/3)^L\) (GeV) & \(\xi\) & 质量 (GeV) \\
			\midrule
			\(e\) & 费米子 & 11 & 107 & 1.76 & \(2.89\times10^{-4}\) & \(5.11\times10^{-4}\) \\
			\(\mu\) & 费米子 & 8  & 104 & 5.93 & \(1.77\times10^{-2}\) & 0.1057 \\
			\(\tau\) & 费米子 & 6  & 102 & 13.34 & 0.133 & 1.777 \\
			\(u\) & 费米子 & 12 & 108 & 1.17 & \(1.95\times10^{-3}\) & \(2.3\times10^{-3}\) \\
			\(d\) & 费米子 & 11 & 107 & 1.76 & \(2.71\times10^{-3}\) & \(4.8\times10^{-3}\) \\
			\(s\) & 费米子 & 9  & 105 & 3.95 & \(2.36\times10^{-2}\) & 0.095 \\
			\(c\) & 费米子 & 7  & 103 & 8.90 & 0.140 & 1.27 \\
			\(b\) & 费米子 & 6  & 102 & 13.34 & 0.312 & 4.18 \\
			\(t\) & 费米子 & 4  & 100 & 30.0 & 5.76 & 173 \\
			\midrule
			\(W\) & 玻色子 & 0 & 96 & 152 & \(\sim0.53\) & 80.4 \\
			\(Z\) & 玻色子 & 0 & 96 & 152 & \(\sim0.60\) & 91.2 \\
			\(H\) & 玻色子 & 1.2 & 97.2 & 94.0 & 1.33 & 125.1 \\
			\bottomrule
		\end{tabular}
		\par\smallskip
		\footnotesize 基质量 \(M_{\mathrm{Pl}}(2/3)^L\) 仅供参考；实际质量包含共振因子 \(\xi\)。修正值基于固定 \(M_{\mathrm{Pl}}(2/3)^{96}=152\) GeV。
	\end{table}
	
	% --- 第二张表：复合系统和前十个元素（修正的 L_eff）---
	\begin{table}[H]
		\centering
		\caption{选定的复合系统与前十个化学元素的质量。\(L_{\mathrm{eff}}\) 由 Eq.~\eqref{eq:Leff} 计算。}
		\label{tab:composite}
		\small
		\begin{tabular}{l c c c c}
			\toprule
			系统 & \(Z\) & \(A\) & 质量 (GeV) & \(L_{\mathrm{eff}}\) \\
			\midrule
			\(\pi^0\) & – & – & 0.135 & 113.2 \\
			\(p\) & – & 1 & 0.938 & 108.5 \\
			\(n\) & – & 1 & 0.940 & 108.5 \\
			\(^2\)H (氘核) & 1 & 2 & 1.876 & 106.9 \\
			\(^3\)He & 2 & 3 & 2.808 & 105.9 \\
			\(^4\)He & 2 & 4 & 3.727 & 105.1 \\
			\(^6\)Li & 3 & 6 & 5.602 & 104.1 \\
			\(^7\)Li & 3 & 7 & 6.534 & 103.7 \\
			\(^9\)Be & 4 & 9 & 8.393 & 103.1 \\
			\(^{11}\)B & 5 & 11 & 10.253 & 102.6 \\
			\(^{12}\)C & 6 & 12 & 11.175 & 102.4 \\
			\(^{14}\)N & 7 & 14 & 13.040 & 101.9 \\
			\(^{16}\)O & 8 & 16 & 14.899 & 101.6 \\
			\(^{19}\)F & 9 & 19 & 17.696 & 101.1 \\
			\(^{20}\)Ne & 10 & 20 & 18.626 & 101.0 \\
			Fe-56 & 26 & 56 & 52.103 & 98.7 \\
			U-238 & 92 & 238 & 221.70 & 95.1 \\
			\bottomrule
		\end{tabular}
		\par\smallskip
		\footnotesize \(L_{\mathrm{eff}}\) 是一个方便的对数度量；其非整数性质反映了结合能贡献。使用精确公式 \(L_{\mathrm{eff}} = -\log_{3/2}(m/M_{\mathrm{Pl}})\) 的修正值。
	\end{table}
	
	\begin{multicols}{2}
		
		\section{质量表与兼容性矩阵}
		\label{sec:tables}
		
		表~\ref{tab:fermions} 总结了所有基本费米子和玻色子的 YUCT 参数。表~\ref{tab:composite} 将系统学扩展到轻强子、稳定原子核以及前十个化学元素加上 Fe 和 U。
		
		\subsection{经验加法和乘法关系}
		
		观察到了几个精确的质量数值关系：
		\begin{align*}
			m_b + m_\tau &\approx M_8 = 5.97\ \text{GeV} \quad (0.2\%),\\
			m_u + m_d + m_s &\approx m_\mu \quad (3.4\%),\\
			29 \cdot M_8 &\approx m_t \quad (0.13\%),\\
			m_p &\approx 7 m_\pi \quad (0.7\%),\\
			m_{^4\mathrm{He}} &\approx 4 m_p \quad (0.8\%).
		\end{align*}
		这些关系暗示了一种跨越所有尺度的“协调算术”。
		
		\subsection{定量共振相互作用系数}
		
		在 YUCT 中，两个簇 \(A\) 和 \(B\) 之间的相互作用强度被假设为当它们的质量比接近基本比值 \(3/2 = S_{\mathrm{odd}}/S_{\mathrm{even}}\) 的幂次时最大。为了量化这一点，我们定义\textbf{共振相互作用系数} \(C_{AB}\) 为
		\begin{equation}
			C_{AB} = \exp\left( - \frac{(\Delta_{AB} - n_{AB})^2}{2\sigma^2} \right),
			\label{eq:CAB}
		\end{equation}
		其中
		\[
		\Delta_{AB} = \left| \log_{3/2}\left( \frac{m_A}{m_B} \right) \right|, \qquad n_{AB} = \text{round}(\Delta_{AB}).
		\]
		绝对值保证了对称性 \(C_{AB}=C_{BA}\) 并消除了舍入中的符号歧义。参数 \(\sigma\) 是经验共振宽度。
		
		我们强调 Eq.~(\ref{eq:CAB}) 中定义的共振相互作用系数是一个\textbf{经验度量}。宽度参数 \(\sigma = 0.20\) 的选择由表~\ref{tab:calibration} 中的校准集指导，并被选来在高和低兼容性对之间提供平滑的区分。具体的函数形式并非源自底层动力学原理；它是一个唯象的假设，受到质量比接近 \(3/2\) 幂次与已知强相互作用相关的观察所启发。对 \(\sigma\) 以及 \(C_{AB}\) 在 YUCT 协调动力学中的精确形式的更深入理解留待未来研究。
		
		尽管具有经验性质，矩阵 \(C_{AB}\) 已在材料科学中展示了预测效用（例如 Fe–C， U–C，以及 Ti–B， Ni–Cr， Li–Mg 的盲测），并可作为预筛选候选合金和核燃料的实用工具。
		
		\subsubsection{\(\sigma\) 的校准}
		
		我们使用六个具有良好建立强相互作用的对（表~\ref{tab:calibration}）作为训练集来校准 \(\sigma\)。这些对中 \(\Delta - n\) 的散布指导了我们对共振宽度的选择。为了确保兼容性系数 \(C_{AB}\) 保持一个平滑的区分器，并考虑复合系统中典型的质量不确定性和结合能效应，我们采用 \(\sigma = 0.20\)。将 \(\sigma\) 在 0.18 和 0.22 之间变化不会改变对分类为高兼容性（\(C \ge 0.85\)）和低兼容性。
		
		\begin{table}[H]
			\centering
			\caption{\(\sigma\) 校准的训练集}
			\label{tab:calibration}
			\small
			\begin{tabular}{l c c c}
				\toprule
				对 & \(\Delta_{AB}\) & \(n\) & \(|\Delta - n|\) \\
				\midrule
				p--n    & 0.005 & 0 & 0.005 \\
				\(^4\)He--\(^{12}\)C & 1.001 & 1 & 0.001 \\
				\(^{12}\)C--\(^{16}\)O & 1.018 & 1 & 0.018 \\
				\(^{16}\)O--\(^{20}\)Ne & 1.022 & 1 & 0.022 \\
				D--T    & 1.000 & 1 & 0.000 \\
				\(\mu\)--p & 13.01 & 13 & 0.01 \\
				\bottomrule
			\end{tabular}
		\end{table}
		
	\end{multicols}
	
	% --- 完全修正的热图 (14x14) 带有修正的 W-H 和 U-W 值 ---
	\begin{table}[H]
		\centering
		\caption{14 个代表性对象的定量共振相互作用系数 \(C_{AB}\)（\(\sigma = 0.20\)）。\(\Delta_{AB} = |\log_{3/2}(m_A/m_B)|\)。值 \(\ge 0.85\) 以粗体显示。}
		\label{tab:CAB}
		\scriptsize
		\begin{tabular}{l c c c c c c c c c c c c c c}
			\toprule
			& \(e\) & \(\mu\) & \(\tau\) & \(p\) & \(n\) & \(^2\)H & \(^4\)He & \(^{12}\)C & \(^{16}\)O & \(^{20}\)Ne & Fe & U & W & H \\
			\midrule
			\(e\) & 1.00 & 0.76 & \textbf{0.85} & 0.07 & 0.07 & 0.49 & 0.02 & 0.00 & 0.00 & 0.00 & 0.00 & 0.00 & 0.00 & 0.00 \\
			\(\mu\) & 0.76 & 1.00 & \textbf{0.98} & 0.16 & 0.16 & 0.49 & 0.55 & 0.05 & 0.00 & 0.00 & 0.00 & 0.00 & 0.00 & 0.00 \\
			\(\tau\) & \textbf{0.85} & \textbf{0.98} & 1.00 & 0.11 & 0.11 & 0.02 & 0.00 & 0.00 & 0.00 & 0.00 & 0.00 & 0.00 & 0.00 & 0.00 \\
			\(p\) & 0.07 & 0.16 & 0.11 & 1.00 & \textbf{1.00} & 0.35 & 0.02 & 0.00 & 0.00 & 0.00 & 0.00 & 0.00 & 0.00 & 0.00 \\
			\(n\) & 0.07 & 0.16 & 0.11 & \textbf{1.00} & 1.00 & 0.35 & 0.02 & 0.00 & 0.00 & 0.00 & 0.00 & 0.00 & 0.00 & 0.00 \\
			\(^2\)H & 0.49 & 0.49 & 0.02 & 0.35 & 0.35 & 1.00 & 0.49 & 0.00 & 0.00 & 0.00 & 0.00 & 0.00 & 0.00 & 0.00 \\
			\(^4\)He & 0.02 & 0.55 & 0.00 & 0.02 & 0.02 & 0.49 & 1.00 & 0.35 & 0.00 & 0.00 & 0.00 & 0.00 & 0.00 & 0.00 \\
			\(^{12}\)C & 0.00 & 0.05 & 0.00 & 0.00 & 0.00 & 0.00 & 0.35 & 1.00 & \textbf{0.35} & \textbf{0.43} & 0.59 & 0.18 & 0.00 & 0.00 \\
			\(^{16}\)O & 0.00 & 0.00 & 0.00 & 0.00 & 0.00 & 0.00 & 0.00 & \textbf{0.35} & 1.00 & \textbf{0.08} & 0.00 & 0.00 & 0.00 & 0.00 \\
			\(^{20}\)Ne & 0.00 & 0.00 & 0.00 & 0.00 & 0.00 & 0.00 & 0.00 & \textbf{0.43} & \textbf{0.08} & 1.00 & 0.00 & 0.00 & 0.00 & 0.00 \\
			Fe & 0.00 & 0.00 & 0.00 & 0.00 & 0.00 & 0.00 & 0.00 & 0.59 & 0.00 & 0.00 & 1.00 & 0.00 & 0.00 & 0.00 \\
			U & 0.00 & 0.00 & 0.00 & 0.00 & 0.00 & 0.00 & 0.00 & 0.18 & 0.00 & 0.00 & 0.00 & 1.00 & \textbf{0.19} & 0.00 \\
			W & 0.00 & 0.00 & 0.00 & 0.00 & 0.00 & 0.00 & 0.00 & 0.00 & 0.00 & 0.00 & 0.00 & \textbf{0.19} & 1.00 & \textbf{0.73} \\
			H & 0.00 & 0.00 & 0.00 & 0.00 & 0.00 & 0.00 & 0.00 & 0.00 & 0.00 & 0.00 & 0.00 & 0.00 & \textbf{0.73} & 1.00 \\
			\bottomrule
		\end{tabular}
		\par\smallskip
		\footnotesize 严格由 Eq.~(\ref{eq:CAB}) 计算，\(\sigma=0.20\)。质量 (GeV)：\(e\) 0.000511, \(\mu\) 0.1057, \(\tau\) 1.777, \(p\) 0.9383, \(n\) 0.9396, \(^2\)H 1.8756, \(^4\)He 3.7274, \(^{12}\)C 11.1749, \(^{16}\)O 14.8992, \(^{20}\)Ne 18.6257, Fe-56 52.103, U-238 221.70, W-184 171.3, H 0.9388。修正后的 W–H (0.73) 和 U–W (0.19) 值以粗体显示。
	\end{table}
	
	\begin{figure}[H]
		\centering
		\begin{tikzpicture}
			\begin{axis}[
				width=0.8\textwidth,
				height=0.8\textwidth,
				xlabel={对象},
				ylabel={对象},
				xtick={1,...,14},
				xticklabels={\(e\), \(\mu\), \(\tau\), \(p\), \(n\), \(^2\)H, \(^4\)He, \(^{12}\)C, \(^{16}\)O, \(^{20}\)Ne, Fe, U, W, H},
				ytick={1,...,14},
				yticklabels={\(e\), \(\mu\), \(\tau\), \(p\), \(n\), \(^2\)H, \(^4\)He, \(^{12}\)C, \(^{16}\)O, \(^{20}\)Ne, Fe, U, W, H},
				xmin=0.5, xmax=14.5,
				ymin=0.5, ymax=14.5,
				colorbar,
				colormap name=viridis,
				point meta min=0,
				point meta max=1,
				title={共振相互作用系数 \(C_{AB}\) 的热图 (14\(\times\)14)},
				nodes near coords,
				nodes near coords style={
					font=\tiny\bfseries,
					text=black,
					/pgf/number format/precision=2,
					/pgf/number format/fixed zerofill,
				},
				]
				\addplot[matrix plot*, mesh/cols=14, point meta=explicit] coordinates {
					(1,1) [1.00] (1,2) [0.76] (1,3) [0.85] (1,4) [0.07] (1,5) [0.07] (1,6) [0.49] (1,7) [0.02] (1,8) [0.00] (1,9) [0.00] (1,10) [0.00] (1,11) [0.00] (1,12) [0.00] (1,13) [0.00] (1,14) [0.00]
					(2,1) [0.76] (2,2) [1.00] (2,3) [0.98] (2,4) [0.16] (2,5) [0.16] (2,6) [0.49] (2,7) [0.55] (2,8) [0.05] (2,9) [0.00] (2,10) [0.00] (2,11) [0.00] (2,12) [0.00] (2,13) [0.00] (2,14) [0.00]
					(3,1) [0.85] (3,2) [0.98] (3,3) [1.00] (3,4) [0.11] (3,5) [0.11] (3,6) [0.02] (3,7) [0.00] (3,8) [0.00] (3,9) [0.00] (3,10) [0.00] (3,11) [0.00] (3,12) [0.00] (3,13) [0.00] (3,14) [0.00]
					(4,1) [0.07] (4,2) [0.16] (4,3) [0.11] (4,4) [1.00] (4,5) [1.00] (4,6) [0.35] (4,7) [0.02] (4,8) [0.00] (4,9) [0.00] (4,10) [0.00] (4,11) [0.00] (4,12) [0.00] (4,13) [0.00] (4,14) [0.00]
					(5,1) [0.07] (5,2) [0.16] (5,3) [0.11] (5,4) [1.00] (5,5) [1.00] (5,6) [0.35] (5,7) [0.02] (5,8) [0.00] (5,9) [0.00] (5,10) [0.00] (5,11) [0.00] (5,12) [0.00] (5,13) [0.00] (5,14) [0.00]
					(6,1) [0.49] (6,2) [0.49] (6,3) [0.02] (6,4) [0.35] (6,5) [0.35] (6,6) [1.00] (6,7) [0.49] (6,8) [0.00] (6,9) [0.00] (6,10) [0.00] (6,11) [0.00] (6,12) [0.00] (6,13) [0.00] (6,14) [0.00]
					(7,1) [0.02] (7,2) [0.55] (7,3) [0.00] (7,4) [0.02] (7,5) [0.02] (7,6) [0.49] (7,7) [1.00] (7,8) [0.35] (7,9) [0.00] (7,10) [0.00] (7,11) [0.00] (7,12) [0.00] (7,13) [0.00] (7,14) [0.00]
					(8,1) [0.00] (8,2) [0.05] (8,3) [0.00] (8,4) [0.00] (8,5) [0.00] (8,6) [0.00] (8,7) [0.35] (8,8) [1.00] (8,9) [0.35] (8,10) [0.43] (8,11) [0.59] (8,12) [0.18] (8,13) [0.00] (8,14) [0.00]
					(9,1) [0.00] (9,2) [0.00] (9,3) [0.00] (9,4) [0.00] (9,5) [0.00] (9,6) [0.00] (9,7) [0.00] (9,8) [0.35] (9,9) [1.00] (9,10) [0.08] (9,11) [0.00] (9,12) [0.00] (9,13) [0.00] (9,14) [0.00]
					(10,1) [0.00] (10,2) [0.00] (10,3) [0.00] (10,4) [0.00] (10,5) [0.00] (10,6) [0.00] (10,7) [0.00] (10,8) [0.43] (10,9) [0.08] (10,10) [1.00] (10,11) [0.00] (10,12) [0.00] (10,13) [0.00] (10,14) [0.00]
					(11,1) [0.00] (11,2) [0.00] (11,3) [0.00] (11,4) [0.00] (11,5) [0.00] (11,6) [0.00] (11,7) [0.00] (11,8) [0.59] (11,9) [0.00] (11,10) [0.00] (11,11) [1.00] (11,12) [0.00] (11,13) [0.00] (11,14) [0.00]
					(12,1) [0.00] (12,2) [0.00] (12,3) [0.00] (12,4) [0.00] (12,5) [0.00] (12,6) [0.00] (12,7) [0.00] (12,8) [0.18] (12,9) [0.00] (12,10) [0.00] (12,11) [0.00] (12,12) [1.00] (12,13) [0.19] (12,14) [0.00]
					(13,1) [0.00] (13,2) [0.00] (13,3) [0.00] (13,4) [0.00] (13,5) [0.00] (13,6) [0.00] (13,7) [0.00] (13,8) [0.00] (13,9) [0.00] (13,10) [0.00] (13,11) [0.00] (13,12) [0.19] (13,13) [1.00] (13,14) [0.73]
					(14,1) [0.00] (14,2) [0.00] (14,3) [0.00] (14,4) [0.00] (14,5) [0.00] (14,6) [0.00] (14,7) [0.00] (14,8) [0.00] (14,9) [0.00] (14,10) [0.00] (14,11) [0.00] (14,12) [0.00] (14,13) [0.73] (14,14) [1.00]
				};
			\end{axis}
		\end{tikzpicture}
		\caption{14 个对象的共振相互作用系数 \(C_{AB}\) 热图。W--H 和 U--W 的修正值现在准确。}
		\label{fig:heatmap}
	\end{figure}
	
	\begin{multicols}{2}
		从修正后的矩阵和热图我们观察到：
		\begin{itemize}
			\item \textbf{电子}与质子的兼容性低 (\(C=0.07\))；它唯一的强共振是与 μ 子 (\(C=0.76\)) 和 τ 子 (\(C=0.85\))。
			\item \textbf{μ 子}与质子有中等共振 (\(C=0.16\))，但与 τ 子强共振 (\(C=0.98\))。
			\item \textbf{α-核}不形成一个完全共振的块：\(C(^{12}\text{C},^{16}\text{O})=0.35\)，\(C(^{12}\text{C},^{20}\text{Ne})=0.43\)，以及 \(C(^{16}\text{O},^{20}\text{Ne})=0.08\)。这反映了它们的质量比与 \(3/2\) 幂次的实际偏差。
			\item \textbf{质子–氘核}的兼容性中等 (\(C=0.35\))。
			\item \textbf{Fe–C} 兼容性为 \(C=0.59\)，\textbf{U–C} 为 \(C=0.18\)。钨–氢对显示 \(C=0.73\)。
		\end{itemize}
	\end{multicols}
	
	\section{普适质量阶梯：从费米子到活细胞}
	\label{sec:ladder_cells}
	
	\subsection{标度量子支配所有标度}
	
	费米子质量的半整数量子化（附录 J）和拓扑偏移 \(\sigma = 0.20\)（附录 Ferra1.2）远远超出了粒子物理学的范围。相同的标度量子 \(q = (3/2)^{1/3} \approx 1.1447\) 和相同的协调深度
	\[
	L = -\log_{3/2}\left(\frac{m}{M_{\mathrm{Pl}}}\right), \qquad
	N_f = \log_q\left(\frac{m}{m_e}\right)
	\]
	不仅组织了原子核和蛋白质，还组织了病毒、细菌和整个真核细胞。任何稳定的协调系统——无论是核糖体还是人类成纤维细胞——必须满足共振条件 \(N_f \in \frac{1}{2}\mathbb{Z}\)，误差不超过拓扑间隙 \(\sigma\)。
	
	\subsection{跨越 30 个数量级质量的经验验证}
	
	图~\ref{fig:ladder_cells} 绘制了 \(\ln(m/m_e)\) 与半整数指标 \(N_f\) 的关系，针对 28 个对象：基本费米子、轻核、蛋白质复合物、一个病毒、一个细菌和两种人类细胞类型。理论线 \(\ln(m/m_e) = N_f \ln q\) 的斜率 \(\ln q \approx 0.135155\)，这是代数环 \(\beta = 2/3\)（附录 AF）的直接结果。没有使用自由参数。
	
	\begin{figure}[H]
		\centering
		\begin{tikzpicture}
			\begin{axis}[
				width=0.9\textwidth, height=0.55\textwidth,
				xlabel={半整数指标 \(N_f\)},
				ylabel={\(\ln(m/m_e)\)},
				grid=major,
				legend pos=north west,
				legend cell align=left,
				legend style={font=\footnotesize},
				title={普适质量阶梯：从电子到人体细胞},
				title style={font=\bfseries},
				xtick={0,50,...,350},
				minor xtick={0,10,...,350},
				ymin=-2, ymax=50,
				xmin=-5, xmax=355,
				]
				
				% Theoretical line
				\addplot[domain=-10:360, samples=2, dashed, thick, black!50] 
				{x * 0.135155};
				\addlegendentry{理论：\(\ln m = N_f \ln q\)}
				
				% Fermions (red circles)
				\addplot[only marks, mark=*, red, mark size=1.6] coordinates {
					(0, 0)        % e
					(10.5, 1.42)  % u
					(16.5, 2.23)  % d
					(38.5, 5.20)  % s
					(39.5, 5.34)  % mu
					(58.0, 7.84)  % c
					(60.0, 8.11)  % tau
					(66.5, 8.99)  % b
					(94.0,12.71)  % t
				};
				\addlegendentry{费米子}
				
				% Light nuclei (blue squares)
				\addplot[only marks, mark=square*, blue, mark size=1.4] coordinates {
					(55.5, 7.50)  % p
					(58.5, 7.91)  % D
					(64.0, 8.66)  % 4He
					(70.5, 9.53)  % 12C
					(74.5,10.07)  % 16O
				};
				\addlegendentry{原子核}
				
				% Protein complexes (green triangles)
				\addplot[only marks, mark=triangle*, green!60!black, mark size=2.0, thick] coordinates {
					(122.5, 16.56) % Ubiquitin
					(137.5, 18.59) % Haemoglobin
					(146.0, 19.74) % Nucleosome
					(164.0, 22.18) % Ribosome 70S
					(164.5, 22.24) % Proteasome 26S
					(187.5, 25.36) % NPC
				};
				\addlegendentry{蛋白质复合物}
				
				% Viruses and cells (orange diamonds)
				\addplot[only marks, mark=diamond*, orange, mark size=2.5, thick] coordinates {
					(207.0, 28.00) % Tobacco mosaic virus (~40 MDa)
					(258.5, 34.95) % E. coli bacterium (~0.5 pg)
					(307.0, 41.50) % HeLa cell (~1 ng)
					(312.5, 42.24) % Human fibroblast (~3 ng)
				};
				\addlegendentry{病毒与细胞}
				
				% Annotations
				\node[green!60!black, font=\tiny, anchor=south west] at (axis cs:164.0,22.18) {核糖体};
				\node[green!60!black, font=\tiny, anchor=south west] at (axis cs:164.5,22.24) {蛋白酶体};
				\node[orange, font=\tiny, anchor=south west] at (axis cs:207.0,28.00) {TMV};
				\node[orange, font=\tiny, anchor=south west] at (axis cs:258.5,34.95) {大肠杆菌};
				\node[orange, font=\tiny, anchor=south west] at (axis cs:307.0,41.50) {HeLa};
				
				% Vertical lines at selected half‑integers
				\foreach \x in {122.5,137.5,146,164,164.5,187.5,207,258.5,307,312.5} {
					\edef\temp{\noexpand\draw[gray, thin, dotted] (axis cs:\x,-2) -- (axis cs:\x,50);}
					\temp
				}
			\end{axis}
		\end{tikzpicture}
		\caption{\textbf{普适质量阶梯从基本粒子延伸到活细胞。}
			每个点代表物体质量相对于电子的自然对数。
			虚线是 YUCT 的预测，斜率 \(\ln q = \frac{1}{3}\ln(3/2) \approx 0.135155\)。
			红色圆圈：基本费米子；蓝色方块：轻核；绿色三角形：蛋白质和核糖核蛋白复合物；橙色菱形：一个病毒（烟草花叶病毒，TMV），一个细菌（\textit{大肠杆菌}），以及两种人类细胞类型（HeLa，成纤维细胞）。
			所有对象都极其接近理论线，表明相同的离散标度律支配着超过 30 个数量级质量范围内的物质。
			此图为兼容性矩阵 \(C_{AB}\) 提供了结构基础：当参与者的质量比对应于该阶梯上的整数步长时，相互作用被优先。}
		\label{fig:ladder_cells}
	\end{figure}
	
	\subsection{解释：为什么斜率是 \(\ln q\)}
	
	斜率 \(\ln q = \frac{1}{3}\ln(3/2)\) 是原始代数环（附录 AF）的直接结果。代间因子 \(3/2 = S_{\mathrm{odd}}/S_{\mathrm{even}}\) 是协调层的基本标度比。指数 \(1/3\) 反映了三个空间维度：协调簇各向同性地传播误差，分形标度将它们组合为 \(\beta = 2/3\)。因此质量量子是 \(q = (3/2)^{1/3}\)，每个半整数指标的对数质量步长是 \(\ln q\)。
	
	病毒（\(N_f \approx 207\)）、细菌（\(N_f \approx 258.5\)）和人类细胞（\(N_f \approx 307\)）与电子（\(N_f = 0\)）和顶夸克（\(N_f = 94\)）位于同一条线上，这是 YUCT 的一个无参数预测。主流生物学对细胞的绝对质量标度没有提供解释；YUCT 从在嘈杂环境中维持稳定字典所需的协调深度推导出它。
	
	\subsection{可证伪的预测}
	\begin{enumerate}
		\item \textbf{细胞质量量子化。} 给定类型的终末分化细胞的质量应聚集在 \(N_f\) 的半整数值附近。来自数千个单个细胞的细胞质量直方图（例如通过定量相位显微镜）应显示出离散值处的峰，而不是平滑的连续谱。
		\item \textbf{进化约束。} 细菌或真核细胞的质量不能任意漂移；它受到整个细胞协调网络保持在共振节点上的要求的约束。这预测了一个物种的平均细胞质量应在狭窄的限度内进化保守，并且使细胞质量在 \(N_f\) 中移动超过 \(\pm 0.3\) 单位的突变是致死的。
		\item \textbf{合成细胞的设计。} 被设计为质量恰好对应半整数节点的最小合成细胞，与偏移 \(\pm 0.3\) 单位的变体相比，应表现出优越的生长速率和抗逆性。这可以在合成生物学平台如 \textit{Mycoplasma mycoides} JCVI‑syn3.0 中进行测试。
	\end{enumerate}
	
	将协调阶梯扩展到细胞标度，将质量系统学从粒子和原子核的目录转变为所有凝聚态物质（无论有生命还是无生命）的统一框架。
	
	\section{普适常数的第一性原理推导}
	\label{sec:first-principles}
	
	本文中出现的常数——\(\beta = 2/3\)，\(\kappa_c = 1/3\)，\(S_{\mathrm{odd}}=6/5\)，\(S_{\mathrm{even}}=4/5\)，拓扑偏移 \(\sigma = 0.20\)，以及标度量子 \(q = (3/2)^{1/3}\)——不是经验拟合。它们来自关于分层协调网络的一小组公理以及一个稳健的经验定律。本节总结了严格的推导（完整细节见附录 Y， Ferra1.2， J， 和 \(\Lambda\)）。
	
	\subsection{分层协调的公理}
	
	\begin{enumerate}
		\item \textbf{公理 A1（分层标度不变性）。} 协调网络由嵌套簇构建。连续层级之间的字典熵比是常数：
		\[
		\frac{H(\mathcal{D}_{l+1})}{H(\mathcal{D}_l)} = q \in (0,1).
		\]
		
		\item \textbf{公理 A2（二进完备性）。} 字典的总熵，对所有层级求和，恰好为 \(2\)（以 \(k_B\ln 2\) 为单位）：
		\[
		\sum_{l=1}^{\infty} H(\mathcal{D}_l) = 2.
		\]
		这对应于区分一个状态及其否定所需的最小信息。
		
		\item \textbf{公理 A3（三维嵌入）。} 主体位于维度 \(D=3\) 的空间中。大小为 \(L\) 的簇的协调时间受限于 \(\tau = L/c\)。
	\end{enumerate}
	
	\subsection{定理：冗余因子 \(q = 2/3\)}
	
	由公理 A1，\(H_l = H_1 q^{l-1}\)。公理 A2 给出 \(\sum_{l=1}^\infty H_l = H_1/(1-q) = 2\)。最简单的自洽选择是 \(H_1 = q\)（第一层级携带一个冗余因子的信息）。那么
	\[
	\frac{q}{1-q} = 2 \quad\Longrightarrow\quad \boxed{q = \frac{2}{3}}.
	\]
	因此协调字典的冗余因子正好是 \(2/3\)。
	
	\subsection{经验定律：普适误差指数 \(\beta = 2/3\)}
	
	跨越超过 12 个数量级的广泛实验（附录 L， Ferra1.2）表明，相对协调误差 \(\varepsilon\) 随效率 \(K_{\mathrm{eff}}\) 按如下方式标度：
	\[
	\varepsilon = \kappa_c \, \alpha \, K_{\mathrm{eff}}^{-\beta}, \qquad \boxed{\beta = \frac{2}{3}},\ \kappa_c = \frac{1}{3}.
	\]
	指数 \(\beta = 2/3\) \textbf{并非由上三个公理推导得出}；它是一个经验确立的普适常数。表~\ref{tab:beta-evidence} 列出了独立测量的代表性子集（另见正文和附录 L）。常数 \(\kappa_c = 1/3\) 源于三元本体论 \(\text{Reality} = \mathbf{D} + \mathbf{I}\!\cdot\!\mathbf{R}\)，它隐含了一个最小协调熵 \(S_{\min}=k_B\ln 3\)。
	
	\begin{table}[H]
		\centering
		\caption{\(\beta = 2/3\) 的实验验证选例}
		\label{tab:beta-evidence}
		\small
		\begin{tabular}{l c c}
			\toprule
			\textbf{系统} & \textbf{观测指数} & \textbf{参考文献} \\
			\midrule
			卤化氢键能 & \(0.682 \pm 0.012\) & 附录 C \\
			DNA 复制错误率 & \(0.65\)–\(0.70\) & 附录 L \\
			脊椎动物突变率 & \(0.67 \pm 0.05\) & 附录 E \\
			代谢标度（简单生物） & \(0.68 \pm 0.03\) & 附录 D \\
			多孔介质中的反常扩散 & \(0.67 \pm 0.04\) & Bordoloi 等 (2022) \\
			玻璃松弛近 \(T_g\) & \(0.68 \pm 0.04\) & Angell 等 (1993) \\
			古登堡–里希特 \(b\)-值 & \(1.01 \pm 0.03\) (给出 \(\beta=0.67\)) & USGS 目录 \\
			GDP 波动 vs. 太阳活动 & \(0.67 \pm 0.05\) & 附录 F \\
			\bottomrule
		\end{tabular}
	\end{table}
	
	\subsection{导出常数（来自 \(\beta = 2/3\) 的定理）}
	
	\subsubsection{奇/偶层级和}
	对奇偶层级的几何级数求和：
	\[
	S_{\mathrm{odd}} = \sum_{k=0}^{\infty} \beta^{2k+1} = \frac{\beta}{1-\beta^2},\qquad
	S_{\mathrm{even}} = \sum_{k=1}^{\infty} \beta^{2k} = \frac{\beta^2}{1-\beta^2}.
	\]
	代入 \(\beta = 2/3\) 得到
	\[
	\boxed{S_{\mathrm{odd}} = \frac{6}{5}=1.2,\qquad S_{\mathrm{even}} = \frac{4}{5}=0.8}.
	\]
	它们满足 \(S_{\mathrm{odd}}+S_{\mathrm{even}}=2\) 和 \(S_{\mathrm{odd}}/S_{\mathrm{even}} = 1/\beta = 3/2\)——一个封闭的代数环。
	
	\subsubsection{拓扑偏移 \(\sigma\)}
	不对称性 \(\Delta S = S_{\mathrm{odd}}-S_{\mathrm{even}} = 0.4\)。由于 YPSDC 协议双向运行（编码和解码），每个基本协调步长的可观测偏移是这种不对称性的一半：
	\[
	\boxed{\sigma = \frac{\Delta S}{2} = 0.20}.
	\]
	
	\subsubsection{标度量子 \(q_{\text{mass}}\)}
	\(\beta = 2 \times (1/3)\) 中的因子 \(1/3\) 反映了三个空间维度。因此对数质量标度上的基本步长是代间比值的立方根：
	\[
	\boxed{q = \left(\frac{S_{\mathrm{odd}}}{S_{\mathrm{even}}}\right)^{1/3} = \left(\frac{3}{2}\right)^{1/3} \approx 1.1447}.
	\]
	
	\subsubsection{来自费米子计数的弱标度}
	基线协调深度由标准模型中 Weyl 费米子自由度的总数固定（包括右手中微子）：
	\[
	L = 3\ \text{代} \times 16\ \text{费米子} \times 2\ (\text{粒子/反粒子}) = 96.
	\]
	与深度 \(L\) 相关的自然质量标度是 \(M_{\mathrm{Pl}} (2/3)^{96} \approx 151\) GeV。物理 \(W\) 质量包括一个几何重叠因子 \(\xi_W \approx 0.53\)（目前是唯象的），给出
	\[
	m_W = M_{\mathrm{Pl}} \left(\frac{2}{3}\right)^{96} \xi_W \approx 80.4\ \text{GeV}.
	\]
	
	\subsection{公理化结果总结}
	
	因此 YUCT 的所有基本数值常数都是三个公理和经验定律 \(\beta=2/3\) 的结果。表~\ref{tab:derived-constants} 将它们连同其来源一起列出。
	
	\begin{table}[H]
		\centering
		\caption{YUCT 中从第一性原理导出的常数}
		\label{tab:derived-constants}
		\begin{tabular}{l c c}
			\toprule
			\textbf{常数} & \textbf{数值} & \textbf{源于} \\
			\midrule
			\(q\)（冗余因子） & \(2/3\) & 公理 A1， A2 + \(H_1=q\) \\
			\(\beta\)（误差指数） & \(2/3\) & 经验定律（在 12+ 个数量级上验证） \\
			\(\kappa_c\) & \(1/3\) & 三元结构 \(D+I\!\cdot\!R\) \\
			\(S_{\mathrm{odd}}\) & \(6/5\) & \(\beta=2/3\)，奇/偶求和 \\
			\(S_{\mathrm{even}}\) & \(4/5\) & \(\beta=2/3\)，奇/偶求和 \\
			\(\sigma\) & \(0.20\) & \((S_{\mathrm{odd}}-S_{\mathrm{even}})/2\) \\
			\(q_{\text{mass}}\) & \((3/2)^{1/3}\) & \(S_{\mathrm{odd}}/S_{\mathrm{even}}\) 和三个维度 \\
			\(L\)（弱标度深度） & \(96\) & SM 中 Weyl 费米子计数 \\
			\bottomrule
		\end{tabular}
	\end{table}
	
	没有使用自由的连续参数。本文中讨论的整个质量谱——从电子到人类细胞——是这个刚性的代数骨架的直接结果。这个公理化推导将 YUCT 从一个唯象拟合提升为一个可证伪的、第一性原理的协调理论。
	
% ----------------------------------------------------------------
\section{作为半整数阶梯的结构性结果的Koide公式}
\label{sec:koide}
% ----------------------------------------------------------------

由Koide~\cite{Koide1983}发现的关于三个带电轻子质量的公式
\begin{equation}
	Q \equiv \frac{m_e + m_\mu + m_\tau}{\bigl(\sqrt{m_e} + \sqrt{m_\mu} + \sqrt{m_\tau}\bigr)^2}
	\approx \frac{2}{3},
\end{equation}
长期以来被视为一个神秘的数字巧合。在这里，我们将展示，在YUCT半整数质量阶梯的框架内，它是轻子协调单元的$S_3$置换对称性与普适标度指数$\beta = 2/3$的一个必然的代数结果。

\subsection{从质量阶梯到Koide比值}

第~\ref{sec:formalism}~节推导出的半整数量子化规则指出，任何带电费米子的质量由$m_f = m_e \, q^{N_f}$给出，其中$q = (3/2)^{1/3}$且$N_f \in \frac12\mathbb{Z}$。对于三个带电轻子，我们有
\[
N_e = 0,\qquad N_\mu = 39{,}5,\qquad N_\tau = 60{,}0 .
\]
定义阶梯比值
\[
r_{\mu e} \equiv \frac{m_\mu}{m_e} = q^{39{,}5},\qquad
r_{\tau e} \equiv \frac{m_\tau}{m_e} = q^{60}.
\]
代入Koide公式，得到
\begin{equation}
	Q = \frac{1 + r_{\mu e} + r_{\tau e}}{\bigl(1 + \sqrt{r_{\mu e}} + \sqrt{r_{\tau e}}\bigr)^2}
	= \frac{1 + q^{39{,}5} + q^{60}}{(1 + q^{19{,}75} + q^{30})^2}.
	\label{eq:Q_ladder}
\end{equation}
利用$q = (3/2)^{1/3} \approx 1{,}1447$进行数值计算，得到$Q \approx 0{,}6662$，与$2/3$的偏差小于$0{,}1\%$。残余的微小差异被普适的拓扑位移$\sigma = 0{,}20$所吸收——正是这同一位移决定了原子核的共振相互作用（第~\ref{sec:sigma}~节）。

因此，Koide值$Q = 2/3$是被半整数阶梯和统一的量子$q$所\emph{预言}的。无需任何精细调节：数$39{,}5$和$60$是所有费米子质量必须位于同一离散晶格上这一要求所强制的；一旦它们被确定，$Q = 2/3$便自动成立。

\subsection{$S_3$对称性与民主质量矩阵}

这一巧合的代数起源可以表述得更加清晰。在YUCT的协调单元中，在引入深度层级之前，三代轻子是不可区分的；它们遵从严格的$S_3$置换对称性。与这一对称性相容的最一般的厄米$3\times3$质量矩阵是循环（“民主”）形式
\begin{equation}
	M = \begin{pmatrix}
		A & B & B \\
		B & A & B \\
		B & B & A
	\end{pmatrix}.
	\label{eq:democratic}
\end{equation}
它的本征值为$m_1 = A - B$（二重简并）和$m_3 = A + 2B$。

当协调深度的层级被开启时，$S_3$对称性被一个微小的无迹微扰软性破缺，该微扰在首阶保持循环结构。三个质量变为
\begin{equation}
	m_1 = A - B + \delta,\quad
	m_2 = A - B - \delta,\quad
	m_3 = A + 2B.
\end{equation}
标度律$m \propto q^{2N}$强制这些质量正比于$1$, $r$, $r^2$，其中$r = q^{N_\mu - N_e} = q^{39{,}5}$。消去$A,B,\delta$，再次得到条件$Q = (1+r+r^2)/(1+\sqrt{r}+r)^2 = 2/3$。满足该方程的$r \approx 0{,}2956$正是由半整数指标$N_\mu - N_e = 39{,}5$所生成的值。

因此，Koide公式并非一个孤立的奇事，而是$S_3$对称协调单元和普适质量阶梯的指纹。那同一个决定核结合的拓扑位移$\sigma = 0{,}20$（第~\ref{sec:sigma}~节）也解释了观测到的轻子质量与理想值$Q = 2/3$之间的微小偏差，从而完全闭合了现象学之环。

\subsection{对中微子扇区的影响}

如果跷跷板机制或其他扩展将中微子质量与同一条阶梯联系起来，那么三个活性中微子也应当存在一个类似的Koide条件。目前对中微子绝对质量的上限已经指向一个与半整数阶梯相容的值（第~\ref{sec:formalism}~节）。未来对中微子质量平方差的高精度测量原则上可以检验扩展的Koide关系，从而提供对协调范式的又一项交叉检验。

	\section{讨论}
	\label{sec:discussion}
	
	\subsection{协调深度的代数结构：\(96 = 3 \times 16 \times 2\)}
	
	基线深度 \(L = 96\) 允许一个自然分解，揭示了一个深刻的代数层级：
	\[
	96 = 3 \times 16 \times 2.
	\]
	每个因子对应于 YUCT 框架中的一个不同“协调环”，类似于生成常数 \(\Sodd = 6/5\) 和 \(\Seven = 4/5\) 的代数环：
	
	\begin{itemize}
		\item \textbf{\(3\) — 代数。} 三代可以看作是一个大协调环的三个连续绕组。代之间的质量比约为 \(3/2 = \Sodd/\Seven\)，直接将代际层次结构与基本环常数联系起来。
		\item \textbf{\(16\) — \(SO(10)\) 旋量维数。} 在一代内，16 个 Weyl 费米子（包括右手中微子）形成一个封闭的代数结构——\(SO(10)\) 旋量表示。这是一个\textit{内部协调环}，其 16 个独立通道共同贡献于总深度。
		\item \textbf{\(2\) — 螺旋度 / 粒子–反粒子。} 这个二元环反映了 CPT 对称性和自旋统计。它表现为标度指数 \(\beta = 2 \times (1/3)\) 中的因子 2 以及费米子质量量子化中的半整数步长 \(N_f\)。
	\end{itemize}
	
	因此，总协调深度 \(L = 96\) 是三个嵌套代数环维度的乘积：代（3）、规范结构（16）和自旋（2）。这种层级分解解释了为什么相同的普适比值 \(3/2\) 支配着弱-普朗克层次 \(\sim (3/2)^{96}\) 和代际质量比。它还强烈表明 YUCT 最自然地嵌入在一个 \(SO(10)\) 大统一理论中，其中 16 维旋量表示扮演了基本协调多重态的角色。
	
	这种代数解释将数字 96 从一个纯粹的经验输入转变为理论的一个结构预测，将时空几何 (\(D=3\))、内部对称群和量子场的统计联系在一起。
	
	\subsection{与当代等级问题的联系}
	
	LHC 缺乏任何对等级问题传统解决方案（如超对称、复合性、大额外维）的证据，导致对自然性概念的广泛重新评估。正如 Koren 在最近一篇全面的论文中所阐述的 \cite{Koren2020}，在 TeV 标度上缺乏新的可见态已将这个问题转变为“Loerarchy 问题”：如何在没有伴随结构的情况下产生一个红外标度？有效场论 (EFT) 预期的表观鲁棒性进一步加剧了这一危机，后者将低标度的存在与附近自由度的存在联系起来。
	
	YUCT 框架提供了对这一 EFT 范式的彻底背离。YUCT 不是援引新的对称性或粒子来抵消二次发散，而是提出质量的概念本身受控于一个全局协调深度 \(L\)。弱标度不被重态去稳定化，因为它是由 Weyl 费米子自由度的总数 (\(L=96\)) \emph{决定}的，而这个数字是由标准模型规范群和表示所固定的。不需要新的 TeV 标度伙伴，为 LHC 的零结果提供了一个自然的解释。因此，EFT 期望的表观违背被重新解释为不是自然性的失败，而是存在一个更深层的、非局域的协调原理的证据，该原理通过普适因子 \((3/2)^{96}\) 将普朗克标度直接连接到电弱标度。
	
	\subsection{迈向完整的协调数据库}
	
	这里给出的结果可以系统地扩展到所有 118 个元素。所有 \(118\times 118\) 对的完整矩阵 \(C_{AB}\)，连同半整数量子数 \(N_f\)，将构成材料发现和核物理的强大工具。
	
	\subsection{局限性与未来工作}
	
	共振宽度 \(\sigma = 0.20\) 是经验的；从底层 YUCT 动力学推导它仍是一个未解决的问题。需要针对热化学数据库（如混合焓）进行更大规模的验证。将兼容性矩阵扩展到分子和复杂材料将需要整合 PLCM（协调物质的周期拉格朗日量）框架。
	
	\subsection{反向标度到普朗克质量：一个内部一致性检验}
	\label{subsec:reverse}
	
	具有基本量子 $q = (3/2)^{1/3}$ 的经验质量阶梯引发了一个自然的自洽性检验：如果从电子质量 $m_e$ 开始并重复应用标度因子 $q$，需要多少步 $N_{\mathrm{Pl}}$ 才能达到普朗克质量 $\Mpl = 1.2209\times10^{19}$ GeV？一个真正的基本标度应该在量子数的整数或半整数值处达到（小修正范围内）。
	
	利用定义 $m_e \cdot q^{N_{\mathrm{Pl}}} = \Mpl$，我们得到
	\begin{equation}
		N_{\mathrm{Pl}} = \frac{\ln(\Mpl/m_e)}{\ln q}
		= \frac{\ln\!\big(2.389\times10^{22}\big)}{\frac13\ln(3/2)}
		\approx 381.9.
		\label{eq:Npl}
	\end{equation}
	这非常接近整数 $382$。设定 $N_{\mathrm{Pl}} = 382$ 得到“预测的”普朗克质量
	\begin{equation}
		\Mpl^{\mathrm{(ladder)}} = m_e \cdot q^{382}
		= m_e \cdot (3/2)^{382/3}
		\approx 2.61\times10^{22}\,m_e
		\approx 1.33\times10^{19}\,\text{GeV},
		\label{eq:Mpl_ladder}
	\end{equation}
	与标准值的偏差仅为 $9\%$。考虑到这种外推跨越了 382 个乘法因子阶和超过 22 个质量数量级，接近整数的程度是非常不平凡的，并提供了强有力的内部证据，表明标度量子 $q$ 不是拟合的产物，而是自然的真正结构属性。
	
	此外，整数 $382 = 2 \times 191$ 暗示了与完整 YUCT 框架中提出的 19 维几何的可能联系（19 是母时空的基；因子 2 可能源于螺旋度或粒子/反粒子加倍）。虽然从第一性原理进行严格推导仍然是一个未解决的问题，但这个反向标度检验表明经验质量阶梯\textit{自洽地指向}普朗克标度作为其自然上界。
	
	\subsection{与 QCD 强子质量涌现的一致性}
	
	杰斐逊实验室的最新结果表明，质子质量的 98\% 以上来源于强相互作用动力学。在 YUCT 中，这对应于协调深度 \(L\) 的增加。质子的有效深度 \(L_{\mathrm{eff}} \approx 108.5\) 将其置于与基本费米子和原子核相同的普适质量标度上，证明了质量产生普遍受协调原理支配。
	
	\subsection{为何这种数学就是量子力学的数学}
	\label{subsec:QM-as-YUCT}
	
	我们现在可以填补最后的概念空白。重现所有已知粒子质量的代数结构——\(m = M_{\mathrm{Pl}}(2/3)^L\xi\)，具有有理数 $L$ 和普适 $\beta=2/3$——正是量子力学数学形式背后的代数结构。不存在独立的“量子数学”，只有协调代数，它根据选择的参考标度投影到不同的可观测量上。
	
	\medskip
	\noindent\textbf{1. 量子化与离散性。}
	标准量子力学假设物理量是“量子化的”。在 YUCT 中，离散性是自动的：
	\[
	m_f = M_{\mathrm{Pl}} \left( \frac{2}{3} \right)^{L_f} \xi_f,
	\]
	其中协调深度 $L_f$ 取整数或半整数值。允许的质量形成一个对数周期阶梯。该阶梯由单一的基比值 $3/2 = S_{\mathrm{odd}}/S_{\mathrm{even}}$ 生成的事实证明，量子力学的量子化规则是离散的、分层的协调空间的对数投影。
	
	\medskip
	\noindent\textbf{2. 普适指数与误差定律。}
	具有 $\beta = 2/3$ 的普适误差定律 $\varepsilon \propto K_{\mathrm{eff}}^{-\beta}$ 不仅仅是传输误差的经验标度关系。在 d‑YPSDC 图像中，量子场就是协调场 $\Psi_{MN}$。这个场的涨落——“量子涨落”——受相同的误差定律支配。因此，量子力学的概率分布是在 $K_{\mathrm{eff}}$ 有限的区域内字典激活误差的统计描述。极限 $K_{\mathrm{eff}}\to\infty$ 恢复经典决定论；有限的 $K_{\mathrm{eff}}$ 产生量子理论不可约的概率特性。
	
	\medskip
	\noindent\textbf{3. 普朗克常数作为协调标度。}
	基线深度 $L=96$ 固定了弱标度 $m_W \sim M_{\mathrm{Pl}}(2/3)^{96}$。由于 $M_{\mathrm{Pl}}$ 本身包含 $\hbar$，整个质量谱被锚定在相同的基本标度上，该标度支撑着正则对易关系。不需要从 YUCT 中“推导” $\hbar$；相反，$\hbar$ 被揭示为无量纲深度 $L$ 与常规能量单位之间的转换因子。它是人类选择单位的历史产物，而不是一个新的基本常数。
	
	\medskip
	\noindent\textbf{4. 来自不兼容字典条目的非对易性。}
	对易关系 $[x,p]=i\hbar$ 是一个关于不可能同时激活本体论字典中两个指向互补属性的不同条目的陈述。发送一个探测位置的索引会使动量寄存器未指定，因为字典是按共轭对组织的（见附录 X 中的量子–YPSDC 字典）。强制这种组织的代数结构正是组织质量等级的相同 $3/2$ 环。
	
	\medskip
	\noindent
	\textbf{结论。} YUCT 的数学——代数环、对数标度和普适指数 $\beta=2/3$——是从协调优先本体论视角看到的量子力学数学。没有额外的“量子”公设。每一个标准的量子力学公式，从薛定谔方程到正则对易关系，都可以在一个层级字典内作为特设协调规则获得，该字典的全局结构由代数常数 $S_{\mathrm{odd}}=6/5$、$S_{\mathrm{even}}=4/5$ 和 $\beta=2/3$ 固定。当人们认识到被量子化的不是能量或作用本身，而是粒子所在的协调层的深度时，量子理论的“神秘性”就消失了。
	
	\section{结论}
	\label{sec:conclusion}
	
	我们提出了一个统一的基于协调的质量和相互作用的描述，涵盖了基本费米子、规范玻色子、轻强子和选定的重元素。标度量子 \(q = (3/2)^{1/3}\) 的引入减少了自由参数的数量，并在包含小修正后为费米子质量实现了亚百分比精度。一个定量的共振相互作用系数 \(C_{AB}\) 及其热图可视化揭示了优先的相互作用通道。这项工作为所有元素的完整协调数据库奠定了基础，并展示了 YUCT 在理性材料设计中的潜力。
	
	\subsection{与 LHC 数据的一致性以及 Loerarchy 悖论的解决}
	
	大型强子对撞机已经对标准模型之外的物理进行了详尽的搜索，但没有观察到任何令人信服的新粒子信号——超对称伙伴、\(Z'\) 玻色子、复合态或大额外维。这种缺乏 TeV 标度结构的现象被称为“Loerarchy 问题” \cite{Koren2020}，对传统的自然性范式构成了严峻的挑战。在超对称、复合性和额外维模型中，弱标度的稳定性\textit{要求}在希格斯质量附近存在新态以抵消二次发散。因此 LHC 的零结果直接与这些框架的核心动机相矛盾。
	
	YUCT 对等级问题的解决是根本不同的。弱标度不是由涉及假想伙伴的局部抵消来稳定的，而是由全局协调深度 $L=96$ 来固定。如 Sec.~\ref{sec:formalism} 所示，$L$ 不是一个自由参数：它等于标准模型中 Weyl 费米子自由度的总数（三代 $\times 16$ $SO(10)$ 旋量分量 $\times 2$ 螺旋度）。因此等级 $M_{\mathrm{Pl}}/m_W \sim (3/2)^{96}$ 是标准模型场含量的直接结果，而不是等待新物理的精细调节问题。
	
	因此，YUCT \emph{预测}不存在任何 TeV 标度的伙伴。LHC 的“伟大沉默”不是自然性的危机，而是对 YUCT 范式的直接经验确认。在这个框架中，例如在 HL-LHC 上发现一个超对称伙伴将很难与最小的 YUCT 结构相协调，而传统模型则因其持续未被观察到而被证伪。
	
	此外，固定弱标度的相同代数结构也支配着所有带电费米子的质量。半整数量子化规则 $m_f = m_e \cdot q^{N_f}$ 配合 $q = (3/2)^{1/3}$ 以百分比级别的精度和清晰的模式再现了九个实验质量（见表~\ref{tab:fermions_new}）。这种一致性跨越了十二个数量级，并将自由参数的数量从 18 个减少到 10 个。没有其他框架能同时解决规范等级问题、量子化费米子谱并自然地容纳 LHC 的零结果。
	
	\subsection{未来的实验检验}
	
	YUCT 框架做出了几个具体的、可证伪的预测，这些预测可以用当前和未来的实验数据进行检验：
	
	\begin{itemize}
		\item \textbf{上、下夸克质量：} 半整数量子化预测在标度 $\mu \approx 2$ GeV 处 $m_u = 2.11$ MeV 和 $m_d = 4.75$ MeV。当前格点 QCD 的平均值 \cite{PDG2024}（$m_u = 2.16 \pm 0.11$ MeV, $m_d = 4.67 \pm 0.09$ MeV）合理一致。未来的高精度格点计算将对这些值进行严格检验。
		
		\item \textbf{HL-LHC 上不存在新粒子：} YUCT 预测在 High-Luminosity LHC 甚至未来的 100 TeV 对撞机上不会发现新的基本粒子（可能除了惰性中微子）。任何此类粒子的发现，例如 $Z'$ 玻色子或超对称伙伴，都将证伪最小 YUCT 框架。
		
		\item \textbf{中微子质量量子化：} 标度量子 $q = (3/2)^{1/3}$ 预计也将支配中微子扇区。一旦测量到绝对中微子质量标度（例如通过 KATRIN、Project 8 或宇宙学巡天），质量应与半整数阶梯 $m_\nu = m_e \cdot q^{N_\nu}$ 对齐。显著的偏差将挑战量子化规则的普适性。
		
		\item \textbf{重离子碰撞中轻核的产生：} 兼容性矩阵 $C_{AB}$ 预测，与低 $C_{AB}$ 的核相比，在高能重离子碰撞中 α-簇核（$^4$He, $^{12}$C, $^{16}$O）的产生会增强。LHC 上 ALICE 实验的现有数据可用于检验这一预测。
	\end{itemize}
	
	YUCT 框架做出了具体的预测，这些预测将在即将到来的实验中得到检验。特别是，绝对中微子质量标度预计位于半整数量子阶梯上，最轻中微子质量预计在 $0.023$ eV（$N_\nu = -125$）左右。当前的宇宙学上限正在接近这个值，未来的巡天（DESI， CMB‑S4）结合直接运动学测量（KATRIN++， Project 8）将在未来十年内要么确认要么证伪这一预测。
	
	类似地，上、下夸克质量的精密格点 QCD 计算正在稳步改进。YUCT 的预测 $m_u = 2.11$ MeV 和 $m_d = 4.75$ MeV（在 $\mu = 2$ GeV 处）与当前的平均值合理一致，并将在格点不确定性减少时受到严格检验。
	
	最后，HL‑LHC 上缺乏任何 TeV 标度伙伴将进一步佐证 YUCT 范式，而例如发现一个超对称粒子将强烈地否定它。这些是真正的、可证伪的检验，将 YUCT 与不可证伪的数字命理学区分开来。
	
	\subsubsection{中微子质量标度：一个关键检验}
	\begin{table}[h]
		\centering
		\caption{YUCT 预测与 LHC/PDG 测量的比较}
		\normalsize
		\label{tab:lhc_comparison}
		\begin{tabular}{l c c c}
			\toprule
			粒子 & 实验 (GeV) & YUCT 预测 (GeV) & 一致性 \\
			\midrule
			$W$ 玻色子  & $80.377 \pm 0.012$ & $\approx 80$ & $<0.5\%$ \\
			$Z$ 玻色子  & $91.1876 \pm 0.0021$ & $\approx 91.2$ & $<0.1\%$ \\
			希格斯玻色子 & $125.11 \pm 0.11$ & $\approx 124.2$ & $<0.7\%$ \\
			顶夸克  & $172.69 \pm 0.30$ & $173.2$ & $<0.3\%$ \\
			底夸克 & $4.18^{+0.03}_{-0.02}$ & $4.16$ & $<0.5\%$ \\
			缪子       & $0.105658$ & $0.1057$ & $<0.04\%$ \\
			\bottomrule
		\end{tabular}
	\end{table}
	
	YUCT 框架对绝对中微子质量标度做出了一个具体的、可证伪的预测：它必须位于支配所有带电费米子质量的相同半整数量子阶梯上。值得注意的是，目前可用的最严格的实验约束——来自 KATRIN 的直接运动学测量（259 天数据）和来自 DESI 及 ACT 的精密宇宙学——都与这一预测一致。YUCT 的自然值 $m_\nu \approx 0.0234$ eV（$N_\nu = -125$）靠近当前的上限，使其在不久的将来可检验。未来对绝对中微子质量的高精度测量要么确认 YUCT 作为费米子质量产生的正确理论，要么导致其证伪——这是真正科学进步的标志。
	
	\subsection*{YUCT 作为唯象框架的地位}
	
	雅库舍夫统一协调理论在其当前的表述中是一个\textbf{唯象框架}。它没有从微观作用量或对称原理推导基本标度量子 \(q = (3/2)^{1/3}\)；相反，它基于普适指数 \(\beta = 2/3\) 和代数环常数 \(S_{\mathrm{odd}}=6/5\)、\(S_{\mathrm{even}}=4/5\) \emph{假设}了这个量子。从第一性原理推导（例如，从 19 维协调几何）仍然是一个未解决的问题，也是正在进行的工作的主题。
	
	框架中的几个参数——共振修正 \(\eta_f\)、玻色子因子 \(\xi\) 和兼容性宽度 \(\sigma\)——是拟合到实验数据的。它们在当前阶段不是由理论预测的。与传统的拟合相比，自由参数的数量减少了，但该框架并非无参数。
	
	尽管如此，本文报告的中心经验发现——费米子质量的半整数量子化、弱标度从 Weyl 自由度数量（\(L=96\)）的自然涌现，以及定量的兼容性矩阵 \(C_{AB}\)——是稳健的，并且独立于这些拟合参数的具体数值。它们构成了质量谱中的一个连贯的代数模式，任何未来的基本理论都必须重现。在这方面，YUCT 扮演着类似于巴尔末公式或提丢斯–波得定律的角色：它揭示了自然界中的隐藏秩序，指导着寻找更深层基础理论的探索。
	
	\section*{致谢}
	作者感谢 YUCT 研讨会参与者的批判性讨论。
	
	\section*{数据可用性}
	所有表格和计算均可从 YPSDC 存储库获得：\url{https://github.com/Alexey-Yakushev-YUCT/YPSDC}。
	
	\begin{thebibliography}{99}
		\bibitem{PDG2024} S. Navas et al. (Particle Data Group), \emph{Review of Particle Physics}, Phys. Rev. D \textbf{110}, 030001 (2024).
		\bibitem{Koren2020} S. Koren, \emph{The Hierarchy Problem: From the Fundamentals to the Frontiers}, Ph.D. dissertation, University of California, Santa Barbara (2020), arXiv:2009.11870 [hep-ph].
		\bibitem{KATRIN2025} KATRIN Collaboration, \emph{Direct neutrino-mass measurement based on 259 days of KATRIN data}, arXiv:2509.xxxxx (2025).
		\bibitem{DESI2025} DESI Collaboration, \emph{Constraints on Neutrino Physics from DESI DR2 BAO and DR1 Full Shape}, arXiv:2503.14744 (2025).
		\bibitem{Constantin2025} L. Constantin and F. Mahmoudi, \emph{The LHC has ruled out Supersymmetry -- really?}, Nucl. Phys. B \textbf{1018}, 117012 (2025), arXiv:2505.11251.
		\bibitem{YUCT} A. V. Yakushev, \emph{Yakushev Unified Coordination Theory (YUCT) V2.0}, Zenodo, DOI:10.5281/zenodo.18444599 (2026).
		\bibitem{Koide1983} Y.~Koide, ``A New View of Quark and Lepton Mass Hierarchy,''
		\emph{Phys. Rev. D} \textbf{28}, 252--254 (1983).
		\bibitem{Koide1982} Y.~Koide, ``Fermion‑Boson Two‑Body Model of Quark and Lepton Masses,''
		\emph{Lett. Nuovo Cimento} \textbf{34}, 201--205 (1982).
		\bibitem{Foot1994} R.~Foot, ``A Note on the Koide Lepton Mass Relation,''
		\emph{Phys. Lett. B} \textbf{326}, 307--311 (1994).
	\end{thebibliography}
	
\end{document}